\section{Implement a Stack using a Linked List}
\label{sec:sq-stack-linked-list}

\problemstatement{Implement a stack supporting $\textit{push}(x)$,
$\textit{pop}()$, $\textit{top}()$, and $\textit{isEmpty}()$, backed by a
singly linked list, each operation running in $O(1)$, with no fixed
capacity limit.}

\begin{patternbox}
The keyword shift from Section~\ref{sec:sq-stack-array} is \emph{``no
fixed capacity''} -- an array-backed stack needs to know its maximum size
up front (or pay for resizing), while a linked list grows and shrinks node
by node with no such limit. The structural insight: a stack's LIFO order
matches a linked list's \textbf{head} perfectly, since inserting and
removing at the head of a singly linked list are both already $O(1)$.
\end{patternbox}

\subsection*{Understanding the problem carefully}

If we always insert new elements at the \textbf{head} of the list and
always remove from the head, the most recently inserted element is always
the first one removed -- exactly LIFO order, with no need to traverse the
list at all.

\begin{ideabox}[Head Insertion = Push, Head Removal = Pop]
Maintain a single pointer \textit{head} to the front of the list (null
when empty). $\textit{push}(x)$ creates a new node with $\textit{next}
\gets \textit{head}$, then sets $\textit{head}$ to this new node.
$\textit{pop}()$ reads $\textit{head}$'s value, then advances
$\textit{head} \gets \textit{head}.\textit{next}$.
\end{ideabox}

\subsection*{Algorithm}

\begin{enumerate}
  \item $\textit{head} \gets$ null initially.
  \item $\textit{push}(x)$: $\textit{node} \gets$ new node with value $x$,
        $\textit{next} \gets \textit{head}$; $\textit{head} \gets
        \textit{node}$.
  \item $\textit{pop}()$: $\textit{val} \gets \textit{head}.\textit{val}$;
        $\textit{head} \gets \textit{head}.\textit{next}$; return
        \textit{val}.
  \item $\textit{top}()$: return $\textit{head}.\textit{val}$.
  \item $\textit{isEmpty}()$: return $(\textit{head} = \text{null})$.
\end{enumerate}

\subsection*{Dry run}

\dryrun{Push $1,2,3$; pop once.}

\begin{itemize}
  \item Push $1$: list $= 1 \to \text{null}$.
  \item Push $2$: list $= 2 \to 1 \to \text{null}$.
  \item Push $3$: list $= 3 \to 2 \to 1 \to \text{null}$.
  \item Pop: returns $3$; list $= 2 \to 1 \to \text{null}$.
\end{itemize}

\subsection*{C++ implementation}

\begin{lstlisting}
#include <bits/stdc++.h>
using namespace std;

class Stack {
    struct Node {
        int val;
        Node* next;
        Node(int v, Node* n) : val(v), next(n) {}
    };
    Node* head = nullptr;

public:
    void push(int x) {
        head = new Node(x, head);
    }

    int pop() {
        int val = head->val;
        Node* old = head;
        head = head->next;
        delete old;
        return val;
    }

    int top() {
        return head->val;
    }

    bool isEmpty() {
        return head == nullptr;
    }
};

int main() {
    Stack s;
    s.push(1); s.push(2); s.push(3);
    cout << "Popped: " << s.pop() << endl;
    cout << "Top: " << s.top() << endl;
    return 0;
}
\end{lstlisting}

\begin{complexitybox}
\textbf{Time Complexity.} $O(1)$ for every operation -- all of them touch
only the head node.
\textbf{Space Complexity.} $O(n)$ for $n$ stored elements, with one node
of constant-size overhead each, and no wasted unused capacity (unlike a
fixed-size array implementation).
\end{complexitybox}

\begin{takeawaybox}
A stack's LIFO discipline matches a singly linked list's head exactly --
no \textit{tail} pointer is ever needed, since insertion and removal both
happen at the same end. This is the cleanest possible linked-list data
structure to implement, and it is exactly why the next section's queue
(needing access at \emph{both} ends) is structurally harder.
\end{takeawaybox}
